\documentclass[runningheads]{llncs}

\usepackage[T1]{fontenc}
\usepackage{lmodern}
\usepackage[utf8]{inputenc}
\usepackage[english,brazilian]{babel}
\usepackage{amsmath,amssymb}
\usepackage{graphicx}
\usepackage{array}
\usepackage{booktabs}
\usepackage{multirow}
\usepackage{tabularx}
\usepackage{url}
\usepackage{tikz}
\usetikzlibrary{arrows.meta,positioning}
\usepackage[round,authoryear]{natbib}
\usepackage[hidelinks]{hyperref}

\renewcommand{\bibsection}{\section*{Referências}}
\renewcommand{\bibfont}{\small}
\setlength{\bibsep}{0pt}
\AtBeginDocument{\providecommand{\doi}{}%
  \renewcommand{\doi}[1]{doi: \href{https://doi.org/#1}{\nolinkurl{#1}}}}
\newcolumntype{L}[1]{>{\raggedright\arraybackslash}p{#1}}

\begin{document}

\title{EMGInvaders Adaptativo: Aprendizado Contínuo de Perfis de Ativação Muscular para Refinar a Calibração em um Jogo de Reabilitação Facial}
\titlerunning{EMGInvaders Adaptativo}

\author{Matheus Kubrusly de Miranda Montenegro da Silva \and Pedro Henrique Paragó Santos
\and Felipe Nogueira Lima Dias}
\authorrunning{~}
\institute{CEFET/RJ}

\maketitle

\begin{abstract}
Jogos controlados por eletromiografia de superfície (sEMG) tornam mais motivadores os exercícios de reabilitação, mas costumam exigir uma calibração a cada sessão, porque o sinal muda com a posição dos eletrodos, com o ambiente em que o jogo é usado e com a fadiga. Guardar a calibração de uma sessão para a outra também não resolve, já que o padrão de contração do paciente muda à medida que ele evolui, e essa evolução precisa ficar registrada. Este trabalho propõe tratar a calibração do EMGInvaders, um jogo controlado por sEMG facial, como um problema de aprendizado contínuo. O perfil de ativação muscular de cada jogador é dividido em duas partes: a faixa do sinal, que vai do nível em repouso ao da contração máxima e é medida de novo a cada sessão, e os parâmetros que decidem qual gesto foi feito, ajustados com poucas gravações de gesto novas e com as gravações guardadas das sessões anteriores. Durante a partida, o limiar de detecção acompanha a fadiga, que fica registrada. A hipótese é que a recolocação dos eletrodos altera sobretudo a faixa do sinal, de modo que as gravações antigas, ajustadas a ela, continuam úteis e permitem uma calibração mais curta e um acerto maior. A avaliação planejada compara a proposta com a versão atual do jogo, em sessões repetidas e em ambientes com mais e menos interferência, quanto ao tempo de preparação, ao acerto dos gestos e à sua consistência entre sessões, aos disparos indevidos, à usabilidade e à capacidade do perfil de manter o que aprendeu nas primeiras sessões. As contribuições esperadas são reduzir o tempo de calibração e refinar a detecção dos gestos de cada jogador.

\keywords{Eletromiografia de superfície \and Reabilitação facial \and Aprendizado contínuo \and Calibração \and Jogos sérios}
\end{abstract}

%=====================================================================
\section{Introdução}\label{sec:introducao}
%=====================================================================

A reabilitação motora depende da repetição frequente de exercícios específicos, e essa repetição tende a reduzir o engajamento do paciente. Jogos sérios, que são jogos com finalidade terapêutica, são usados para tornar essa rotina mais motivadora~\citep{lohse2013video}, principalmente quando mostram ao paciente, em tempo real, um sinal fisiológico medido nele próprio, o que permite que ele aprenda a controlar a função medida. Esse recurso é chamado de \textit{biofeedback}~\citep{giggins2013biofeedback}.

A eletromiografia de superfície (sEMG) registra, por eletrodos colocados sobre a pele, a atividade elétrica que acompanha a contração muscular~\citep{hermens2000seniam}. Por ser não invasiva, ela é adequada ao uso repetido, embora o custo, a pouca facilidade de uso dos equipamentos e o tempo exigido por cada medição ainda limitem sua adoção clínica~\citep{campanini2020barriers}. Na face, ela permite medir de forma objetiva a função de pacientes com lesão do nervo facial~\citep{franz2024facial}. Nesse contexto se insere o EMGInvaders~\citep{sousa2026emginvaders}, um jogo inspirado no \textit{Space Invaders} controlado pela contração da musculatura facial. Na versão atual, dois sensores de sEMG de baixo custo, um de cada lado do rosto, ligam-se a um Arduino: a contração de um lado move a nave para esse lado, e a dos dois lados ao mesmo tempo dispara. Essa versão já funciona nos testes feitos durante o desenvolvimento, mas ainda não foi avaliada com voluntários nem em sessões repetidas, etapa que esta proposta planeja.

A dificuldade desses sistemas está no próprio sinal. A relação entre o esforço do paciente e o valor medido muda de uma sessão para outra, porque pequenos deslocamentos dos eletrodos já alteram as leituras~\citep{young2011electrode}. Por isso, antes de cada partida, o EMGInvaders mede de novo, de cada lado do rosto, o nível do sinal em repouso e na contração máxima, e esse tempo sai do tempo de terapia. Durante a partida, a fadiga altera a amplitude do sinal~\citep{cifrek2009fatigue}, e um limiar fixo deixa de corresponder ao esforço do paciente. A versão atual já reajusta a amplitude de referência pelos picos mais recentes e adapta três parâmetros de decisão às últimas sessões do jogador, mas descarta as sessões mais antigas a cada atualização. Além disso, o jogo pode ser usado em lugares diferentes. Nos testes feitos durante o desenvolvimento, afastar cabos de energia do sensor já mudava o sinal de forma visível, mas o efeito do ambiente ainda não foi medido.

Guardar a calibração e reaproveitá-la na sessão seguinte não resolve o problema, porque o padrão de ativação reflete o estado clínico do paciente e muda com o tratamento~\citep{machetanz2022facial}. Com a calibração congelada, o jogo continua descrevendo o paciente do primeiro dia e não acompanha a recuperação. Se, ao contrário, o jogo se adapta sem guardar nada, a evolução que se quer acompanhar se perde. Esse conflito é estudado pelo aprendizado contínuo como o dilema entre plasticidade, que é a capacidade de incorporar o que é novo, e estabilidade, que é a de preservar o que já foi aprendido~\citep{parisi2019continual}.

Há ainda outra dificuldade. Quando o sinal muda, nem sempre dá para saber se o eletrodo foi recolocado, uma mudança que o jogo deve compensar, ou se o paciente melhorou, uma mudança que deve ficar registrada. Para lidar com isso, a proposta usa dois elementos: o perfil de ativação muscular, ou simplesmente perfil, que descreve como os músculos de cada jogador se ativam em repouso, na contração máxima e com o cansaço ao longo da partida, como isso muda de uma sessão para outra e quais valores o jogo usa para reconhecer os gestos dele, e as gravações de gesto, que são trechos do sinal gravados enquanto o jogo pede um gesto, guardados junto com o nome do gesto pedido. A pergunta de pesquisa é a seguinte: \textit{em sessões repetidas do EMGInvaders, atualizar o perfil com poucas gravações de gesto novas e com as gravações guardadas das sessões anteriores reduz o tempo de preparação e aumenta o acerto em relação à versão atual, sem apagar o histórico do participante?} A hipótese central é que a recolocação dos eletrodos altera sobretudo a faixa do sinal, que vai do nível em repouso ao da contração máxima de cada lado. Se isso for verdade, as gravações antigas, ajustadas à faixa da sessão em que foram feitas, continuam representando o jogador. Reaproveitá-las permite trocar a calibração completa por uma curta e, como o perfil passa a se apoiar em mais gravações do que as de uma única sessão, também aumentar o acerto e torná-lo mais estável de uma sessão para outra.

O objetivo geral é propor e avaliar um mecanismo de aprendizado contínuo do perfil de ativação muscular para o EMGInvaders. Os objetivos específicos são: (i)~separar a faixa do sinal de cada sessão dos parâmetros de decisão do jogador; (ii)~atualizar o perfil a cada sessão com poucas gravações novas e com as guardadas, sem descartar as sessões anteriores; (iii)~ajustar o limiar de ativação à fadiga e registrá-la como dado da sessão; (iv)~registrar as condições de cada ambiente e medir como a interferência afeta a detecção; e (v)~definir um protocolo que compare a proposta com a versão atual. As Seções~\ref{sec:revisao} e~\ref{sec:relacionados} apresentam os conceitos e os trabalhos relacionados; as Seções~\ref{sec:metodologia} e~\ref{sec:avaliacao}, a proposta e sua avaliação; e a Seção~\ref{sec:conclusao}, as contribuições esperadas, as limitações e o cronograma.

%=====================================================================
\section{Revisão da Literatura}\label{sec:revisao}
%=====================================================================

\subsection{Eletromiografia de Superfície e Fontes de Mudança no Sinal}\label{sec:semg}

Na contração, as unidades motoras recrutadas geram potenciais elétricos. Um par de eletrodos sobre o músculo mede a diferença entre eles, e um terceiro, de referência, fica em região eletricamente inativa~\citep{hermens2000seniam}. Para controlar um jogo, costuma-se usar o envelope do sinal, versão retificada e suavizada que representa a intensidade da contração. Na face, os músculos são finos e próximos, o que favorece a captação da atividade de músculos vizinhos, o \textit{crosstalk}, e torna o posicionamento dos eletrodos mais crítico~\citep{franz2024facial}.

Três fontes de mudança alteram o sinal, cada uma com um tratamento diferente. O reposicionamento dos eletrodos entre sessões degrada os sistemas treinados para reconhecer padrões de contração~\citep{young2011electrode}. A fadiga, dentro da sessão, altera a amplitude e o conteúdo de frequência do sinal~\citep{cifrek2009fatigue}. A recuperação motora modifica o próprio padrão de ativação: na paralisia facial, o índice de lateralidade, que compara a amplitude dos dois lados do rosto, acompanha a gravidade clínica~\citep{machetanz2022facial}. O ambiente em que o sistema é usado, com a sua interferência elétrica, não está entre essas fontes, e por isso o protocolo o registra.

Se a recolocação dos eletrodos altera sobretudo a faixa do sinal, ajustar cada sessão à sua própria faixa preserva o restante da informação sobre o modo de contrair do jogador. A literatura consultada não confirma essa distinção para a sEMG facial, e por isso ela é a hipótese que o protocolo da Seção~\ref{sec:avaliacao} testa.

\subsection{Aprendizado Contínuo}\label{sec:aprendizado}

O aprendizado contínuo trata de incorporar informação nova ao longo do tempo sem descartar o que já foi aprendido~\citep{parisi2019continual}. Sua dificuldade central é o esquecimento catastrófico: ao ajustar um modelo aos dados recentes, o desempenho sobre os dados antigos cai de forma abrupta~\citep{kirkpatrick2017overcoming}. As estratégias formam três famílias~\citep{parisi2019continual}: a regularização, que restringe a mudança dos parâmetros importantes para as tarefas antigas~\citep{kirkpatrick2017overcoming}; as arquiteturas dinâmicas, que ampliam o modelo; e o \textit{replay}, que reapresenta exemplos guardados junto aos dados novos. Num perfil de poucos parâmetros, como o do EMGInvaders, a regularização e o \textit{replay} se traduzem em formas simples de atualização, e as arquiteturas dinâmicas, voltadas a redes neurais, não se aplicam.

A questão prática do \textit{replay} é escolher os exemplos que entram na memória, de capacidade limitada. No iCaRL, eles são escolhidos de modo que sua média se aproxime da média da classe~\citep{rebuffi2017icarl}; o sorteio, ao contrário, ignora o conteúdo e pode deixar a memória pouco representativa.

No EMGInvaders, cada sessão corresponde a uma tarefa e cada gesto a uma classe. A plasticidade consiste em acompanhar a mudança do modo de contrair do jogador, e a estabilidade, em preservar o registro das sessões anteriores. O esquecimento aparece quando o perfil deixa de reconhecer as gravações das primeiras sessões, o que a retenção do perfil, definida na Seção~\ref{sec:avaliacao}, mede. A memória de \textit{replay} guarda gravações de cada gesto, e escolhê-las pela média do gesto, como no iCaRL, é uma forma candidata de mantê-la representativa.

%=====================================================================
\section{Trabalhos Relacionados}\label{sec:relacionados}
%=====================================================================

{\sloppy A seleção combina uma busca na base Scopus com a análise das referências do trabalho mais próximo (\textit{snowballing} para trás). As buscas cruzam \texttt{("surface electromyography" OR sEMG)} com \texttt{(game OR "serious game") AND rehabilitation} e com \texttt{("continual learning" OR "incremental learning" OR "transfer learning")}. Incluem-se artigos revisados por pares, em inglês, sobre sEMG em jogos ou \textit{biofeedback} de reabilitação ou sobre a variação do sinal entre sessões ou pessoas.\par}

Os jogos de reabilitação por sEMG tratam cada sessão de forma isolada. \citet{seo2022emg} treinam o padrão de ativação da mão de pacientes após acidente vascular cerebral. A cada sessão, os eletrodos são recolocados e o participante repete um exercício de calibração. \citet{boulay2024serious} combinam um jogo com \textit{feedback} de sEMG e fisioterapia em crianças com paralisia cerebral, ajustando a dificuldade pela contração voluntária máxima medida antes de cada sessão. Na face, \citet{machetanz2022facial} constroem um \textit{biofeedback} com microcontrolador e módulos de sEMG de baixo custo nos músculos frontal e zigomático, o hardware mais próximo do EMGInvaders, mas não tratam da variação entre sessões.

Os trabalhos sobre a variabilidade do sinal concentram-se nos movimentos da mão. \citet{coteallard2019transfer} reaproveitam, por transferência de aprendizado, dados de vários usuários para reduzir a gravação exigida de cada novo usuário, e \citet{chen2024lifelong} aplicam aprendizado contínuo à estimativa da cinemática da mão entre pessoas. O mais próximo desta proposta é \citet{zhang2024replay}, que enfrenta o deslocamento de eletrodos entre dias com \textit{replay} e escolhe as amostras guardadas por agrupamento (DBSCAN), e não por sorteio. Esses trabalhos usam redes neurais profundas sobre sinais do membro superior, e nenhum distingue a mudança causada pelo eletrodo da causada pelo progresso do paciente.

A Tabela~\ref{tab:relacionados} resume a comparação. Nenhum dos trabalhos encontrados trata da sEMG facial num jogo de reabilitação com perfil persistente, nem separa a mudança da faixa do sinal, que deve ser compensada, da evolução do paciente, que deve ser registrada. O \textit{replay} em si já foi aplicado à sEMG por \citet{zhang2024replay}. O que esta proposta acrescenta é o seu uso num jogo facial de baixo custo, com as gravações guardadas ajustadas à faixa de cada sessão e a evolução acompanhada por medidas clínicas.

\begin{table}[htb]
\centering
\caption{Comparação entre os trabalhos relacionados e a proposta, na ordem em que aparecem no texto.}
\label{tab:relacionados}
\scriptsize
\begin{tabularx}{\textwidth}{@{}L{3.1cm}L{1.7cm}L{1.55cm}XL{1.45cm}@{}}
\toprule
\textbf{Trabalho} & \textbf{Região} & \textbf{Uso} & \textbf{Sessões sucessivas} & \textbf{Separa artefato e evolução} \\
\midrule
\citet{seo2022emg} & Mão & Jogo & Calibração a cada sessão & Não \\
\citet{boulay2024serious} & Perna & Jogo & Calibração a cada sessão & Não \\
\citet{machetanz2022facial} & Face & \textit{Biofeedback} & Não tratadas & Não \\
\citet{coteallard2019transfer} & Antebraço & Gestos & Transferência entre pessoas & Não \\
\citet{chen2024lifelong} & Membro sup. & Cinemática & Aprendizado contínuo & Não \\
\citet{zhang2024replay} & Membro sup. & Gestos & \textit{Replay} entre dias & Não \\
\midrule
Esta proposta & Face & Jogo & Perfil atualizado entre sessões & Sim (a avaliar) \\
\bottomrule
\end{tabularx}
\end{table}

%=====================================================================
\section{Proposta de Metodologia}\label{sec:metodologia}
%=====================================================================

\subsection{Visão Geral}

A proposta organiza o EMGInvaders em duas camadas (Figura~\ref{fig:arquitetura}). Na camada de tempo real, os dois sensores, que já entregam o envelope do sinal, enviam os canais ao Arduino, e, a cada quadro do jogo, o detector compara o envelope com o limiar. A montagem prevista posiciona um par de eletrodos sobre o zigomático maior de cada lado do rosto, com um eletrodo de referência em região eletricamente inativa. Esse músculo da mímica, que eleva o canto da boca, é inervado pelo nervo facial, o alvo da reabilitação. Na versão atual, o músculo captado ainda não foi confirmado.

A camada de adaptação funciona fora do laço do jogo. No início da sessão, ela registra a medida de interferência do ambiente, conduz a calibração curta e devolve ao detector o perfil atualizado. Ao final, guarda na memória as gravações de gesto da sessão. O participante entra por um login, que carrega o seu perfil, e os dados ficam vinculados ao login, e não ao nome da pessoa.

\begin{figure}[t]
\centering
\begin{tikzpicture}[
    node distance=0.75cm and 0.5cm,
    box/.style={draw, rounded corners=2pt, minimum width=2.3cm,
                minimum height=0.8cm, align=center, font=\scriptsize},
    mem/.style={box, fill=gray!20},
    ar/.style={-{Latex[length=1.8mm]}, semithick},
    arb/.style={ar, dashed},
    lab/.style={font=\tiny, inner sep=1pt},
    lay/.style={font=\scriptsize\itshape, rotate=90, text=black!65}
]
\node[box] (sensor) {Módulos de sEMG\\envelope do sinal};
\node[box, right=of sensor] (arduino) {Arduino\\digitalização};
\node[box, right=of arduino] (det) {Detecção\\limiar adaptativo};
\node[box, right=of det] (game) {Jogo\\ação da nave};
\node[box, below=of arduino] (verif) {Calibração\\curta};
\node[mem, below=of det] (mem) {Memória\\gravações de gesto};
\node[box, below=of game] (cons) {Atualização\\da memória};

\draw[ar] (sensor) -- (arduino);
\draw[ar] (arduino) -- (det);
\draw[ar] (det) -- (game);
\draw[ar] (arduino) -- node[lab, left] {início} (verif);
\draw[ar] (mem) -- (verif);
\draw[ar] (verif.north east) -- node[lab, right, pos=0.25] {perfil} (det.south west);
\draw[ar] (game) -- node[lab, right] {fim} (cons);
\draw[arb] (cons) -- (mem);

\node[lay] at ([xshift=-0.35cm]sensor.west) {Tempo real};
\node[lay] at ([xshift=-0.35cm]sensor.west |- verif) {Adaptação};
\end{tikzpicture}
\caption{Arquitetura proposta. A seta tracejada indica a atualização da memória ao final da sessão, e o retângulo cinza, o dado guardado entre sessões.}
\label{fig:arquitetura}
\end{figure}

\subsection{Perfil e Atualização entre Sessões}

O perfil tem duas partes, completadas pelo registro da fadiga de cada sessão (Seção~\ref{sec:limiar}). A primeira é a faixa do sinal de cada lado do rosto, que muda com a recolocação dos eletrodos e, por isso, é medida de novo a cada sessão. Como o nível em repouso é medido no próprio local, espera-se que essa medição também acompanhe parte da interferência do ambiente, o que a análise do efeito do ambiente verifica. A segunda são os parâmetros de decisão, como o limiar de ativação e o equilíbrio entre os lados que distingue o tiro do movimento. Eles descrevem o modo de o jogador contrair e são atualizados ao longo das sessões, o que permite ao jogo acompanhar, por exemplo, o fortalecimento da musculatura.

No início de cada sessão, depois do login, uma calibração curta pede poucas gravações de cada gesto. Elas servem para medir a faixa da sessão e, junto com as gravações das sessões anteriores, guardadas já ajustadas à faixa em que foram feitas, para reajustar os parâmetros de decisão. Se o acerto nessas gravações ficar abaixo de um critério, o jogo volta à calibração completa. O número de gravações e o critério de acerto são definidos no estudo piloto. O critério é a escolha mais delicada: se for permissivo, aceita perfis que já não descrevem o jogador e, se for rigoroso, acaba com o ganho de tempo.

A pesquisa compara três formas de atualizar os parâmetros: a janela das sessões mais recentes, usada pela versão atual, que descarta as antigas; o \textit{replay} das gravações guardadas; e a atualização ponderada pelo histórico, que mantém os parâmetros novos próximos dos anteriores, na linha da regularização. Como a memória tem capacidade limitada, também definida no piloto, compara-se ainda a escolha das gravações guardadas: pela proximidade com a média de cada gesto, como no iCaRL, por agrupamento, como \citet{zhang2024replay}, ou por sorteio, como referência. As escolhas usam a retenção do perfil e o acerto, definidos na Seção~\ref{sec:avaliacao}.

Para separar o artefato da evolução clínica, a faixa de cada sessão não entra na memória, porque depende da montagem. Já a trajetória dos parâmetros de decisão e o índice de lateralidade apresentado na Seção~\ref{sec:semg} são registrados a cada sessão como dados de evolução; a sensibilidade do índice à recolocação dos eletrodos é verificada no piloto.

\subsection{Limiar Adaptativo}\label{sec:limiar}

A detecção de ativação por comparação com um limiar acima do nível de repouso é um procedimento clássico na sEMG~\citep{hodges1996onset}, mas um limiar fixo deixa de servir quando a fadiga altera a amplitude das contrações. O ajuste do limiar precisa, então, atender a dois requisitos: acompanhar essa mudança ao longo da partida e mantê-la visível nos dados, porque a fadiga é uma informação que o terapeuta deseja observar.

A pesquisa compara duas formas de ajuste: o reajuste contínuo pela amplitude das contrações recentes, já usado na versão atual, e a recalibração em pausas curtas da partida, com uma contração de referência. Nas duas, o ajuste fica limitado a um intervalo em torno do valor calibrado, para não esconder a fadiga, e a diferença entre os dois valores é registrada como medida da deriva. A escolha ocorre no piloto, pelo critério de manter estável a taxa de gestos reconhecidos sem aumentar os disparos em repouso.

%=====================================================================
\section{Proposta de Avaliação Experimental}\label{sec:avaliacao}
%=====================================================================

\subsection{Delineamento e Aspectos Éticos}

A avaliação compara duas condições com os mesmos participantes. Na condição~A, o \textit{baseline}, o EMGInvaders funciona como na versão atual: calibração completa a cada sessão, como também fazem os jogos de \citet{seo2022emg} e \citet{boulay2024serious}, ajuste às sessões recentes e reajuste da amplitude pelos picos recentes. Na condição~C, o jogo usa o perfil atualizado por aprendizado contínuo, a calibração curta e o limiar adaptativo com registro da deriva.

O delineamento previsto é cruzado: cada participante realiza uma sequência de sessões em cada condição, separadas por dias para que ocorra o reposicionamento dos eletrodos, e a ordem das condições é sorteada para equilibrar o aprendizado do jogo. Ele é preferido a dois grupos independentes porque, com poucos participantes, comparar cada pessoa consigo mesma reduz o peso das diferenças individuais. O desenho é revisto depois do estudo piloto. Na primeira etapa, os participantes são adultos saudáveis, recorte que viabiliza a coleta e a aprovação ética. Na primeira sessão de~C, o perfil é criado pela calibração completa, e as comparações usam as seguintes, com montagem, ajuste de ganho dos sensores e versão, ritmo e duração da partida fixos e documentados. O ambiente, ao contrário, não é fixado: as sessões acontecem em locais com mais e com menos interferência elétrica, cada participante passa pelos mesmos locais nas duas condições, e cada sessão registra o local e uma medida da interferência, como o nível do sinal durante um repouso instruído e uma ficha do local (computador na tomada ou na bateria, rede cabeada e aparelhos ligados por perto). As medidas de interferência candidatas são escolhidas no piloto. O número de participantes e de sessões é definido por análise de poder a partir do estudo piloto, com poder de 80\% e o nível de significância já corrigido para o número de comparações.

Como a coleta envolve seres humanos e dados vinculados ao participante ao longo do tempo, ela depende da aprovação do Comitê de Ética em Pesquisa e da assinatura do Termo de Consentimento Livre e Esclarecido (TCLE), conforme as Resoluções 466/2012 e 510/2016 do Conselho Nacional de Saúde. O protocolo prevê explicitamente as sessões repetidas e o armazenamento do perfil vinculado ao login.

\subsection{Métricas e Análise}

A Tabela~\ref{tab:metricas} lista as variáveis de resposta, que cobrem o custo de preparação, o funcionamento do perfil, a qualidade e a consistência da detecção e a percepção de uso, medida pela escala SUS (\textit{System Usability Scale}), um questionário de dez itens com escore de 0 a 100~\citep{bangor2008sus}. As variáveis de detecção são medidas num modo de treino rotulado, já existente na versão atual, em que o jogo pede os gestos em ordem sorteada e grava o pedido junto com o reconhecido, o que fornece um gabarito.

\begin{table}[htb]
\centering
\caption{Variáveis de resposta e expectativa da condição~C em relação à condição~A.}
\label{tab:metricas}
\scriptsize
\begin{tabularx}{\textwidth}{@{}L{2.3cm}XL{2.3cm}@{}}
\toprule
\textbf{Aspecto} & \textbf{Variável de resposta} & \textbf{Esperado (C/A)} \\
\midrule
Preparação & Tempo do fim da colocação dos eletrodos ao início da partida & Menor \\
Calibração curta & Sessões em que a calibração curta atinge o critério de acerto & Medida só em C \\
Retenção do perfil & Acerto nas gravações da sessão~1 após atualizações sucessivas & Sem queda \\
Acerto & Gestos reconhecidos iguais aos pedidos no treino rotulado & Maior \\
Consistência & Variação do acerto entre as sessões do mesmo participante & Menor \\
Confiabilidade & Ações disparadas durante repouso instruído & Igual ou menor \\
Fadiga & Gestos reconhecidos no último terço da partida / no primeiro & Menor queda \\
Usabilidade & Escore SUS ao fim de cada sessão & Igual ou maior \\
\bottomrule
\end{tabularx}
\end{table}

A retenção do perfil mede se o perfil esquece o que aprendeu, e não o comportamento do jogador. Depois das atualizações pelas sessões~2 a~$n$, o perfil é usado para reconhecer as gravações da sessão~1. Se o acerto nelas cair, o perfil perdeu o modo de contrair registrado no início. Calculada para cada forma de atualização da Seção~\ref{sec:metodologia}, inclusive a janela da versão atual, essa medida mostra qual delas preserva o histórico.

O teste estatístico é escolhido pelo diagnóstico dos dados: se a normalidade das diferenças pareadas entre A e C se sustenta (teste de Shapiro-Wilk e gráfico Q-Q), usa-se o teste $t$ pareado; caso contrário, o de Wilcoxon pareado, ambos bilaterais. Como todas as comparações entre A e C da Tabela~\ref{tab:metricas} são principais, o nível de significância de 0,05 é corrigido pelo método de Holm, que limita a chance de ao menos um falso positivo entre elas. A calibração curta, que só existe em~C, é descrita sem teste. Para entender o efeito do ambiente, o acerto de cada sessão é relacionado à medida de interferência registrada, em cada condição, numa análise exploratória. Cada comparação relata a diferença, o intervalo de confiança de 95\% e um tamanho de efeito, que informa a magnitude da diferença. Esse procedimento e os critérios de exclusão de sessões são definidos antes da coleta, para que a análise não dependa dos dados observados.

Quanto às ameaças à validade interna, o aprendizado do jogo ao longo das sessões é controlado pelo sorteio da ordem das condições, a diferença entre os ambientes, pela passagem de cada participante pelos mesmos locais nas duas condições, e o fato de quem posiciona os eletrodos conhecer a condição, por um procedimento escrito de montagem. Quanto à validade externa, a amostra de adultos saudáveis não permite observar a evolução causada pela recuperação, e os resultados se limitam à montagem, ao hardware e aos ambientes descritos.

%=====================================================================
\section{Considerações Finais e Cronograma}\label{sec:conclusao}
%=====================================================================

Espera-se que o trabalho contribua com a formulação da calibração entre sessões em sEMG como um problema de aprendizado contínuo (Seção~\ref{sec:aprendizado}), com a separação entre a faixa do sinal de cada sessão, tratada como artefato, e os parâmetros e as medidas clínicas, tratados como evolução, e com um protocolo de avaliação replicável, que inclui medidas de retenção do perfil, de consistência entre sessões e do efeito do ambiente.

Algumas limitações são conhecidas desde o planejamento. A hipótese de que a recolocação altera sobretudo a faixa do sinal pode não se confirmar na face, onde o \textit{crosstalk} é frequente. Nesse caso, as gravações guardadas perdem valor, e a proposta precisa de outra forma de ajuste. O aumento do acerto também pode não aparecer, porque uma calibração completa feita na própria sessão pode descrever o jogador melhor do que as gravações antigas. A aquisição se limita ao envelope de dois canais, o que exclui as análises de frequência e impede separar a interferência da rede elétrica do sinal muscular, e a amostra de adultos saudáveis restringe a avaliação à estabilidade e à retenção. Como trabalhos futuros, a coleta com pacientes permitirá avaliar a plasticidade, e a detecção do descolamento de um eletrodo durante a sessão poderá reajustar a faixa do sinal de imediato.

A Tabela~\ref{tab:cronograma} apresenta o cronograma, de janeiro a junho de 2027. A submissão ao Comitê de Ética ocorre no primeiro mês, em paralelo à revisão da literatura, porque é o caminho crítico: sem a aprovação, nem o piloto nem a coleta podem começar.

\begin{table}[htb]
\centering
\caption{Cronograma estratégico do projeto (janeiro a junho de 2027).}
\label{tab:cronograma}
\scriptsize
\begin{tabular}{@{}l c c c c c c@{}}
\toprule
\multirow{2}{*}{\textbf{Atividade}} & \multicolumn{6}{c}{\textbf{Meses}} \\
\cmidrule(l){2-7}
 & jan & fev & mar & abr & mai & jun \\
\midrule
Revisão da literatura e submissão ao Comitê de Ética & X & X & & & & \\
Implementação do perfil e da atualização entre sessões & & X & X & & & \\
Estudo piloto e definição dos parâmetros e das medidas de ambiente & & & X & & & \\
Limiar adaptativo e integração ao jogo & & & X & X & & \\
Coleta em sessões repetidas e em ambientes diferentes & & & & X & X & \\
Análise dos resultados e escrita do artigo & & & & & X & X \\
Reuniões com o orientador & & X & & X & & X \\
\bottomrule
\end{tabular}
\end{table}

\section*{Agradecimentos}

Os autores agradecem à Coordenação de Aperfeiçoamento de Pessoal de Nível Superior (CAPES) pelo acesso ao Portal de Periódicos.

\bibliographystyle{plainnat}
\bibliography{referencias}

\section*{Uso de Inteligência Artificial Generativa}

Os autores utilizaram IA generativa (Claude Code) para construir as buscas na base Scopus e para auxiliar na localização de referências.

\end{document}
